% ═══════════════════════════════════════════════════════════════
% QR-CODES: Simulationen Reduktion
% Chemie Klasse 8 | DS 3
% ═══════════════════════════════════════════════════════════════

\documentclass[a4paper, 11pt]{article}

\usepackage[utf8]{inputenc}
\usepackage[T1]{fontenc}
\usepackage[ngerman]{babel}

\usepackage{helvet}
\renewcommand{\familydefault}{\sfdefault}

\usepackage[margin=1.5cm]{geometry}
\usepackage{qrcode}
\usepackage[hidelinks]{hyperref}
\usepackage{xcolor}
\usepackage{array}
\usepackage{tabularx}

\definecolor{chemie}{HTML}{2a7a4b}

\setlength{\parindent}{0pt}

% URLs definieren
\newcommand{\simAurl}{https://devbox42.github.io/sciencesim/chemie/kl08/03-reduktion/SIM-03a-teilchenebene.html}
\newcommand{\simBurl}{https://devbox42.github.io/sciencesim/chemie/kl08/03-reduktion/SIM-03b-reaktionsgleichung.html}
\newcommand{\simCurl}{https://devbox42.github.io/sciencesim/chemie/kl08/03-reduktion/SIM-03c-vergleich.html}

\begin{document}

\pagestyle{empty}

\begin{center}

{\LARGE\bfseries\textcolor{chemie}{Simulationen: Reduktion}}

\vspace{0.3cm}

{\large Chemie Klasse 8 | DS 3}

\vspace{1cm}

% ═══════════════════════════════════════════════════════════════
% SIM-03a: Teilchenebene
% ═══════════════════════════════════════════════════════════════

\begin{tabular}{>{\centering\arraybackslash}p{5cm} p{10cm}}

\href{\simAurl}{%
    \qrcode[height=4cm]{\simAurl}%
} &

\vspace{0.5cm}
{\Large\bfseries\textcolor{chemie}{SIM-03a: Teilchenebene}}

\vspace{0.3cm}

Was passiert beim Erhitzen von Silberoxid auf atomarer Ebene? Die Simulation zeigt in 5 Schritten, wie das Ionengitter zerfällt und Sauerstoff entweicht.

\vspace{0.5cm}

\href{\simAurl}{\textcolor{chemie}{\texttt{\scriptsize\simAurl}}}

\\

\end{tabular}

\vspace{0.8cm}
\hrule
\vspace{0.8cm}

% ═══════════════════════════════════════════════════════════════
% SIM-03b: Reaktionsgleichung
% ═══════════════════════════════════════════════════════════════

\begin{tabular}{>{\centering\arraybackslash}p{5cm} p{10cm}}

\href{\simBurl}{%
    \qrcode[height=4cm]{\simBurl}%
} &

\vspace{0.5cm}
{\Large\bfseries\textcolor{chemie}{SIM-03b: Reaktionsgleichung}}

\vspace{0.3cm}

Schritt für Schritt die Reaktionsgleichung aufstellen und ausgleichen. Mit Atombilanz-Zähler zur Kontrolle.

\vspace{0.5cm}

\href{\simBurl}{\textcolor{chemie}{\texttt{\scriptsize\simBurl}}}

\\

\end{tabular}

\vspace{0.8cm}
\hrule
\vspace{0.8cm}

% ═══════════════════════════════════════════════════════════════
% SIM-03c: Vergleich
% ═══════════════════════════════════════════════════════════════

\begin{tabular}{>{\centering\arraybackslash}p{5cm} p{10cm}}

\href{\simCurl}{%
    \qrcode[height=4cm]{\simCurl}%
} &

\vspace{0.5cm}
{\Large\bfseries\textcolor{chemie}{SIM-03c: Vergleich Ox/Red}}

\vspace{0.3cm}

Oxidation und Reduktion im direkten Vergleich. Animierte Darstellung der Sauerstoff-Aufnahme bzw. -Abgabe.

\vspace{0.5cm}

\href{\simCurl}{\textcolor{chemie}{\texttt{\scriptsize\simCurl}}}

\\

\end{tabular}

\vspace{1cm}

\hrule

\vspace{0.5cm}

{\small Scanne die QR-Codes oder klicke auf die Links, um die Simulationen zu öffnen.}

\end{center}

\end{document}
